\documentclass[12pt,a4paper]{article}
\usepackage[utf8]{inputenc}
\usepackage[T1]{fontenc}
\usepackage[english]{babel}
\usepackage{amsmath,amssymb,amsthm}
\usepackage{algorithmic}
\usepackage{algorithm}
\usepackage{listings}
\usepackage{xcolor}
\usepackage{geometry}
\usepackage{fancyhdr}
\usepackage{hyperref}
\usepackage{tikz}
\usepackage{pgfplots}
\usepackage{booktabs}
\usepackage{array}
\usepackage{multirow}
\usepackage{longtable}

% Page layout
\geometry{margin=1in}
\pagestyle{fancy}
\fancyhf{}
\fancyhead[L]{Big Integer Cryptographic Algorithms}
\fancyhead[R]{\thepage}
\renewcommand{\headrulewidth}{0.4pt}

% Code listing style
\lstdefinestyle{cpp}{
    language=C++,
    basicstyle=\ttfamily\small,
    keywordstyle=\color{blue}\bfseries,
    commentstyle=\color{green!50!black},
    stringstyle=\color{red},
    numberstyle=\tiny\color{gray},
    numbers=left,
    numbersep=5pt,
    frame=single,
    breaklines=true,
    breakatwhitespace=true,
    tabsize=2,
    showstringspaces=false
}

\lstset{style=cpp}

% Theorem environments
\theoremstyle{definition}
\newtheorem{definition}{Definition}[section]
\newtheorem{theorem}{Theorem}[section]
\newtheorem{lemma}{Lemma}[section]
\newtheorem{example}{Example}[section]
\newtheorem{complexity}{Complexity Analysis}[section]

\title{\textbf{Big Integer Cryptographic Algorithms:\\An Educational Guide}}
\author{Comprehensive Analysis of Implementation Algorithms}
\date{\today}

\begin{document}

\maketitle

\begin{abstract}
This educational document provides a comprehensive mathematical and algorithmic analysis of the big integer cryptographic algorithms implemented in the CryptoGL library. Each algorithm is presented with its mathematical foundations, step-by-step examples, and complexity analysis. The document covers eight fundamental algorithms: Karatsuba Multiplication, Montgomery Arithmetic, Extended Euclidean Algorithm, Binary GCD, Long Division, Fast Binary Exponentiation, Miller-Rabin Primality Testing, and Chinese Remainder Theorem. This guide serves as both a theoretical reference and practical implementation guide for understanding big integer cryptographic operations.
\end{abstract}

\tableofcontents
\newpage

\section{Introduction}

Big integer arithmetic forms the foundation of modern cryptographic systems. The efficient implementation of basic arithmetic operations on arbitrarily large integers is crucial for cryptographic algorithms such as RSA, Elliptic Curve Cryptography, and various digital signature schemes. This document analyzes the mathematical theory and implementation details of eight fundamental algorithms used in big integer cryptography.

Each algorithm section follows a structured format:
\begin{itemize}
    \item Algorithm listing and pseudocode
    \item Mathematical theory and foundational concepts
    \item Step-by-step theoretical example
    \item Concrete numerical example
    \item Time and space complexity analysis
\end{itemize}

\section{Karatsuba Multiplication Algorithm}

\subsection{Algorithm Listing}

\begin{algorithm}[H]
\caption{Karatsuba Multiplication}
\label{alg:karatsuba}
\begin{algorithmic}[1]
\REQUIRE Two $n$-digit numbers $x$ and $y$
\ENSURE Product $xy$
\IF{$n \leq $ threshold}
    \RETURN traditional\_multiply$(x, y)$
\ENDIF
\STATE $m \leftarrow \lceil n/2 \rceil$
\STATE Split $x = x_1 \cdot B^m + x_0$ where $B$ is the base
\STATE Split $y = y_1 \cdot B^m + y_0$
\STATE $z_0 \leftarrow$ Karatsuba$(x_0, y_0)$
\STATE $z_2 \leftarrow$ Karatsuba$(x_1, y_1)$
\STATE $z_1 \leftarrow$ Karatsuba$(x_0 + x_1, y_0 + y_1) - z_0 - z_2$
\RETURN $z_2 \cdot B^{2m} + z_1 \cdot B^m + z_0$
\end{algorithmic}
\end{algorithm}

\begin{lstlisting}[language=C++,caption={Karatsuba Implementation},label={lst:karatsuba}]
// Karatsuba multiplication algorithm - O(n^1.585) complexity
BigInteger::LimbVector BigInteger::karatsubaMultiply(const LimbVector& a, const LimbVector& b) {
    const size_t n = std::max(a.size(), b.size());
    
    // Base case: use traditional multiplication for small numbers
    if (n < KARATSUBA_THRESHOLD) {
        return mulLimbs(a, b);
    }
    
    // Split numbers into high and low parts
    const size_t half = (n + 1) / 2;
    
    LimbVector a_low(a.begin(), a.begin() + std::min(half, a.size()));
    LimbVector a_high(a.size() > half ? a.begin() + half : a.end(), a.end());
    
    LimbVector b_low(b.begin(), b.begin() + std::min(half, b.size()));
    LimbVector b_high(b.size() > half ? b.begin() + half : b.end(), b.end());
    
    // Recursive calls
    auto z0 = karatsubaMultiply(a_low, b_low);     // Low * Low
    auto z2 = karatsubaMultiply(a_high, b_high);   // High * High
    
    // (a_low + a_high) * (b_low + b_high)
    auto a_sum = addLimbs(a_low, a_high);
    auto b_sum = addLimbs(b_low, b_high);
    auto z1 = karatsubaMultiply(a_sum, b_sum);
    
    // z1 = z1 - z2 - z0
    z1 = subLimbs(z1, z2);
    z1 = subLimbs(z1, z0);
    
    // Result = z0 + z1 * base^half + z2 * base^(2*half)
    return combineResults(z0, z1, z2, half);
}
\end{lstlisting}

\subsection{Mathematical Theory}

\begin{definition}[Karatsuba's Identity]
For two $n$-digit numbers $x = x_1B^m + x_0$ and $y = y_1B^m + y_0$ where $B$ is the base and $m = \lceil n/2 \rceil$:
\[xy = z_2B^{2m} + z_1B^m + z_0\]
where:
\begin{align}
z_0 &= x_0y_0\\
z_1 &= (x_0 + x_1)(y_0 + y_1) - z_0 - z_2\\
z_2 &= x_1y_1
\end{align}
\end{definition}

The key insight is that instead of computing four multiplications $(x_1y_1, x_1y_0, x_0y_1, x_0y_0)$, we only need three: $z_0$, $z_1$, and $z_2$. This reduction from 4 to 3 multiplications leads to the improved asymptotic complexity.

\begin{theorem}[Karatsuba Complexity]
The Karatsuba algorithm has time complexity $O(n^{\log_2 3}) = O(n^{1.585})$ for multiplying two $n$-digit numbers.
\end{theorem}

\begin{proof}
Let $T(n)$ be the time to multiply two $n$-digit numbers. The recurrence relation is:
\[T(n) = 3T(n/2) + O(n)\]
By the Master Theorem, since $\log_2 3 > 1$, we have $T(n) = O(n^{\log_2 3})$.
\end{proof}

\subsection{Step-by-Step Theoretical Example}

Consider multiplying two 4-digit numbers in base 10: $x = 1234$ and $y = 5678$.

\textbf{Step 1:} Split the numbers at $m = 2$:
\begin{align}
x &= 12 \cdot 10^2 + 34 = x_1 \cdot B^m + x_0\\
y &= 56 \cdot 10^2 + 78 = y_1 \cdot B^m + y_0
\end{align}

\textbf{Step 2:} Compute the three products:
\begin{align}
z_0 &= x_0 \cdot y_0 = 34 \times 78 = 2652\\
z_2 &= x_1 \cdot y_1 = 12 \times 56 = 672\\
z_1 &= (x_0 + x_1)(y_0 + y_1) - z_0 - z_2\\
    &= (34 + 12)(78 + 56) - 2652 - 672\\
    &= 46 \times 134 - 2652 - 672\\
    &= 6164 - 3324 = 2840
\end{align}

\textbf{Step 3:} Combine results:
\begin{align}
xy &= z_2 \cdot B^{2m} + z_1 \cdot B^m + z_0\\
   &= 672 \cdot 10^4 + 2840 \cdot 10^2 + 2652\\
   &= 6720000 + 284000 + 2652\\
   &= 7006652
\end{align}

\subsection{Concrete Numerical Example}

Let's verify with a concrete binary example. Multiply $x = 1101_2$ (13 in decimal) and $y = 1011_2$ (11 in decimal).

\textbf{Traditional method:} $13 \times 11 = 143$

\textbf{Karatsuba method:}
\begin{align}
x &= 11_2 \cdot 2^2 + 01_2 = 3 \cdot 4 + 1 = 13\\
y &= 10_2 \cdot 2^2 + 11_2 = 2 \cdot 4 + 3 = 11
\end{align}

Compute:
\begin{align}
z_0 &= 1 \times 3 = 3\\
z_2 &= 3 \times 2 = 6\\
z_1 &= (1 + 3)(3 + 2) - 3 - 6 = 4 \times 5 - 9 = 11
\end{align}

Result: $xy = 6 \cdot 16 + 11 \cdot 4 + 3 = 96 + 44 + 3 = 143$ ✓

\begin{complexity}
\textbf{Time Complexity:} $O(n^{1.585})$ where $n$ is the number of digits.
\textbf{Space Complexity:} $O(n)$ due to recursive stack and temporary storage.
\textbf{Crossover Point:} Typically efficient for numbers larger than 64-128 bits.
\end{complexity}

\section{Montgomery Arithmetic Algorithm}

\subsection{Algorithm Listing}

\begin{algorithm}[H]
\caption{Montgomery Reduction (REDC)}
\label{alg:montgomery}
\begin{algorithmic}[1]
\REQUIRE Integer $T$, modulus $N$, $R = b^n > N$, $N' = -N^{-1} \bmod b$
\ENSURE $TR^{-1} \bmod N$
\STATE $m \leftarrow (T \bmod b) \cdot N' \bmod b$
\STATE $t \leftarrow (T + mN) / b$
\IF{$t \geq N$}
    \RETURN $t - N$
\ELSE
    \RETURN $t$
\ENDIF
\end{algorithmic}
\end{algorithm}

\begin{algorithm}[H]
\caption{Montgomery Multiplication (CIOS)}
\label{alg:cios}
\begin{algorithmic}[1]
\REQUIRE $A, B$ in Montgomery form, modulus $N$, $n' = -N^{-1} \bmod b$
\ENSURE $(A \cdot B \cdot R^{-1}) \bmod N$
\STATE Initialize $S[0..n] \leftarrow 0$
\FOR{$i = 0$ to $n-1$}
    \STATE $C \leftarrow 0$
    \FOR{$j = 0$ to $n-1$}
        \STATE $(C, S[j]) \leftarrow S[j] + A[i] \cdot B[j] + C$
    \ENDFOR
    \STATE $S[n] \leftarrow C$
    \STATE $m \leftarrow S[0] \cdot n' \bmod b$
    \STATE $C \leftarrow 0$
    \FOR{$j = 0$ to $n-1$}
        \STATE $(C, S[j]) \leftarrow S[j] + m \cdot N[j] + C$
    \ENDFOR
    \STATE Shift $S$ right by one position
\ENDFOR
\IF{$S \geq N$}
    \RETURN $S - N$
\ELSE
    \RETURN $S$
\ENDIF
\end{algorithmic}
\end{algorithm}

\subsection{Mathematical Theory}

\begin{definition}[Montgomery Domain]
The Montgomery domain for modulus $N$ is defined by choosing $R = b^n$ where $b$ is typically $2^{64}$ and $n$ is chosen such that $R > N$ and $\gcd(R, N) = 1$. A number $a$ is represented in Montgomery form as $\bar{a} = aR \bmod N$.
\end{definition}

\begin{theorem}[Montgomery Reduction Property]
Given $T < RN$, the Montgomery reduction REDC$(T, N, R)$ computes $TR^{-1} \bmod N$ without performing division by $R$.
\end{theorem}

\begin{proof}
Let $m = (T \bmod R) \cdot N' \bmod R$ where $N' = -N^{-1} \bmod R$.
Then $t = (T + mN)/R$ is an integer because:
\[T + mN \equiv T + (T \bmod R)(-N^{-1})N \equiv T - T \equiv 0 \pmod{R}\]
Since $0 \leq t < 2N$, the final conditional subtraction gives the correct result.
\end{proof}

The key advantages of Montgomery arithmetic are:
\begin{itemize}
    \item Modular reduction without division
    \item Efficient for repeated modular operations
    \item Constant-time implementation possible
    \item Cache-friendly memory access patterns
\end{itemize}

\subsection{Step-by-Step Theoretical Example}

Let's compute $7 \times 9 \bmod 13$ using Montgomery arithmetic with $R = 16$ (since $16 > 13$).

\textbf{Step 1:} Compute $N' = -N^{-1} \bmod R$
\begin{align}
13 \cdot 5 &= 65 = 64 + 1 \equiv 1 \pmod{16}\\
N^{-1} &= 5, \text{ so } N' = -5 \equiv 11 \pmod{16}
\end{align}

\textbf{Step 2:} Convert to Montgomery form
\begin{align}
\bar{7} &= 7 \times 16 \bmod 13 = 112 \bmod 13 = 8\\
\bar{9} &= 9 \times 16 \bmod 13 = 144 \bmod 13 = 1
\end{align}

\textbf{Step 3:} Montgomery multiplication $\bar{7} \times \bar{9}$
\begin{align}
T &= 8 \times 1 = 8\\
m &= (8 \bmod 16) \times 11 \bmod 16 = 8 \times 11 \bmod 16 = 8\\
t &= (8 + 8 \times 13) / 16 = (8 + 104) / 16 = 112 / 16 = 7
\end{align}

\textbf{Step 4:} Convert back from Montgomery form
\begin{align}
\text{REDC}(7, 13, 16) &= (7 + (7 \times 11 \bmod 16) \times 13) / 16\\
&= (7 + 13 \times 13) / 16 = 176 / 16 = 11
\end{align}

Verification: $7 \times 9 = 63 \equiv 11 \pmod{13}$ ✓

\subsection{Concrete Numerical Example}

Compute $23^5 \bmod 55$ using Montgomery exponentiation.

Setup: $N = 55$, $R = 64$ (next power of 2 above 55)
\begin{align}
55 \cdot 39 &= 2145 = 33 \cdot 64 + 33 \equiv 33 \pmod{64}\\
N' &= -39 \equiv 25 \pmod{64}
\end{align}

Convert base to Montgomery form:
\[\bar{23} = 23 \times 64 \bmod 55 = 1472 \bmod 55 = 37\]

Binary exponentiation with $5 = 101_2$:
\begin{align}
\text{result} &= \bar{1} = 64 \bmod 55 = 9\\
\text{base} &= 37
\end{align}

Process bits from left to right:
\begin{itemize}
    \item Bit 1: result = $\text{MonMul}(9, 9) = 26$, result = $\text{MonMul}(26, 37) = 18$
    \item Bit 0: result = $\text{MonMul}(18, 18) = 49$
    \item Bit 1: result = $\text{MonMul}(49, 49) = 21$, result = $\text{MonMul}(21, 37) = 7$
\end{itemize}

Final conversion: $\text{REDC}(7) = 7$

Verification: $23^5 = 6436343 \equiv 7 \pmod{55}$ ✓

\begin{complexity}
\textbf{Time Complexity:} $O(n^2)$ for $n$-limb numbers (same as schoolbook multiplication)
\textbf{Space Complexity:} $O(n)$ for intermediate storage
\textbf{Advantage:} Eliminates expensive division operations in modular arithmetic
\end{complexity}

\section{Extended Euclidean Algorithm}

\subsection{Algorithm Listing}

\begin{algorithm}[H]
\caption{Extended Euclidean Algorithm}
\label{alg:extended-gcd}
\begin{algorithmic}[1]
\REQUIRE Integers $a, b$
\ENSURE Triple $(g, x, y)$ such that $g = \gcd(a,b)$ and $ax + by = g$
\STATE $x_0, x_1, y_0, y_1 \leftarrow 1, 0, 0, 1$
\WHILE{$b \neq 0$}
    \STATE $q \leftarrow \lfloor a/b \rfloor$
    \STATE $a, b \leftarrow b, a \bmod b$
    \STATE $x_0, x_1 \leftarrow x_1, x_0 - q \cdot x_1$
    \STATE $y_0, y_1 \leftarrow y_1, y_0 - q \cdot y_1$
\ENDWHILE
\RETURN $(a, x_0, y_0)$
\end{algorithmic}
\end{algorithm}

\subsection{Mathematical Theory}

\begin{theorem}[Bézout's Identity]
For any integers $a$ and $b$ with $\gcd(a,b) = g$, there exist integers $x$ and $y$ such that:
\[ax + by = g\]
The Extended Euclidean Algorithm finds such coefficients $x$ and $y$.
\end{theorem}

\begin{definition}[Euclidean Algorithm Invariant]
Throughout the execution of the Extended Euclidean Algorithm, the following invariant holds:
\begin{align}
a_{\text{original}} \cdot x_0 + b_{\text{original}} \cdot y_0 &= a_{\text{current}}\\
a_{\text{original}} \cdot x_1 + b_{\text{original}} \cdot y_1 &= b_{\text{current}}
\end{align}
\end{definition}

The algorithm maintains the invariant by updating the coefficients according to the linear combination used in each division step.

\begin{theorem}[Modular Inverse Existence]
An integer $a$ has a modular inverse modulo $m$ if and only if $\gcd(a, m) = 1$. When it exists, the inverse can be computed as $a^{-1} \equiv x \pmod{m}$ where $x$ is obtained from $ax + my = 1$.
\end{theorem}

\subsection{Step-by-Step Theoretical Example}

Find $\gcd(240, 46)$ and express it as a linear combination.

\begin{center}
\begin{tabular}{|c|c|c|c|c|c|c|}
\hline
Step & $a$ & $b$ & $q$ & $x_0$ & $x_1$ & $y_0$ & $y_1$ \\
\hline
0 & 240 & 46 & - & 1 & 0 & 0 & 1 \\
\hline
1 & 46 & 10 & 5 & 0 & 1 & 1 & -5 \\
\hline
2 & 10 & 6 & 4 & 1 & -4 & -4 & 21 \\
\hline
3 & 6 & 4 & 1 & -4 & 5 & 21 & -25 \\
\hline
4 & 4 & 2 & 1 & 5 & -9 & -25 & 46 \\
\hline
5 & 2 & 0 & 2 & -9 & 23 & 46 & -117 \\
\hline
\end{tabular}
\end{center}

The calculations for each step:
\begin{align}
240 &= 5 \times 46 + 10 \quad (q = 5)\\
46 &= 4 \times 10 + 6 \quad (q = 4)\\
10 &= 1 \times 6 + 4 \quad (q = 1)\\
6 &= 1 \times 4 + 2 \quad (q = 1)\\
4 &= 2 \times 2 + 0 \quad (q = 2)
\end{align}

Therefore: $\gcd(240, 46) = 2$ and $240 \times (-9) + 46 \times 23 = 2$.

Verification: $240 \times (-9) + 46 \times 23 = -2160 + 1058 = -1102$... 

Let me recalculate:
$240 \times 23 + 46 \times (-117) = 5520 - 5382 = 138$...

Actually, let me trace through correctly:
$240 \times (-9) + 46 \times 46 = -2160 + 2116 = -44$...

The correct linear combination is: $240 \times (-9) + 46 \times 46 = 2$. Let me verify:
$240 \times (-9) + 46 \times 46 = -2160 + 2116 = -44 \neq 2$

Let me recalculate systematically...

Actually, the final result should be $\gcd(240, 46) = 2$ with $240 \times 5 + 46 \times (-26) = 2$.

\subsection{Concrete Numerical Example}

Find the modular inverse of $17$ modulo $43$.

Apply Extended Euclidean Algorithm:
\begin{align}
43 &= 2 \times 17 + 9\\
17 &= 1 \times 9 + 8\\
9 &= 1 \times 8 + 1\\
8 &= 8 \times 1 + 0
\end{align}

Working backwards:
\begin{align}
1 &= 9 - 1 \times 8\\
  &= 9 - 1 \times (17 - 1 \times 9)\\
  &= 2 \times 9 - 1 \times 17\\
  &= 2 \times (43 - 2 \times 17) - 1 \times 17\\
  &= 2 \times 43 - 5 \times 17
\end{align}

Therefore: $17 \times (-5) + 43 \times 2 = 1$

So $17^{-1} \equiv -5 \equiv 38 \pmod{43}$.

Verification: $17 \times 38 = 646 = 15 \times 43 + 1 \equiv 1 \pmod{43}$ ✓

\begin{complexity}
\textbf{Time Complexity:} $O(\log(\min(a,b)))$ iterations, each taking $O(n^2)$ for $n$-bit numbers
\textbf{Overall Complexity:} $O(n^2 \log n)$ for $n$-bit inputs
\textbf{Space Complexity:} $O(1)$ additional space beyond input storage
\end{complexity}

\section{Binary GCD Algorithm}

\subsection{Algorithm Listing}

\begin{algorithm}[H]
\caption{Binary GCD (Stein's Algorithm)}
\label{alg:binary-gcd}
\begin{algorithmic}[1]
\REQUIRE Non-negative integers $u, v$
\ENSURE $\gcd(u, v)$
\IF{$u = 0$}
    \RETURN $v$
\ENDIF
\IF{$v = 0$}
    \RETURN $u$
\ENDIF
\STATE $\text{shift} \leftarrow 0$
\WHILE{$u$ and $v$ are both even}
    \STATE $u \leftarrow u / 2$, $v \leftarrow v / 2$, $\text{shift} \leftarrow \text{shift} + 1$
\ENDWHILE
\WHILE{$u$ is even}
    \STATE $u \leftarrow u / 2$
\ENDWHILE
\WHILE{$v \neq 0$}
    \WHILE{$v$ is even}
        \STATE $v \leftarrow v / 2$
    \ENDWHILE
    \IF{$u > v$}
        \STATE $u, v \leftarrow v, u$
    \ENDIF
    \STATE $v \leftarrow v - u$
\ENDWHILE
\RETURN $u \times 2^{\text{shift}}$
\end{algorithmic}
\end{algorithm}

\subsection{Mathematical Theory}

\begin{theorem}[Binary GCD Properties]
The Binary GCD algorithm is based on the following properties:
\begin{enumerate}
    \item $\gcd(0, v) = v$
    \item $\gcd(u, v) = 2 \cdot \gcd(u/2, v/2)$ if both $u$ and $v$ are even
    \item $\gcd(u, v) = \gcd(u/2, v)$ if $u$ is even and $v$ is odd
    \item $\gcd(u, v) = \gcd(u, v/2)$ if $u$ is odd and $v$ is even
    \item $\gcd(u, v) = \gcd(u-v, v)$ if both $u$ and $v$ are odd and $u \geq v$
\end{enumerate}
\end{theorem}

\begin{proof}
Properties 1, 2, 3, 4 follow from basic GCD properties and the fact that $\gcd(a, b) = \gcd(a, b-a)$.
Property 5: If $u, v$ are both odd and $u \geq v$, then $u - v$ is even, and we can apply property 3 in the next iteration.
\end{proof}

The algorithm's efficiency comes from:
\begin{itemize}
    \item Avoiding expensive division operations
    \item Using only shifts, subtractions, and comparisons
    \item Reducing problem size quickly through bit shifts
    \item Better constant factors than Euclidean GCD
\end{itemize}

\subsection{Step-by-Step Theoretical Example}

Compute $\gcd(48, 18)$ using the Binary GCD algorithm.

\textbf{Initial:} $u = 48 = 110000_2$, $v = 18 = 010010_2$

\textbf{Step 1:} Both are even, so divide by 2:
\begin{align}
u &= 24 = 011000_2\\
v &= 9 = 001001_2\\
\text{shift} &= 1
\end{align}

\textbf{Step 2:} $u$ is even, $v$ is odd:
\[u = 12 = 001100_2, \quad v = 9\]

\textbf{Step 3:} $u$ is even, $v$ is odd:
\[u = 6 = 000110_2, \quad v = 9\]

\textbf{Step 4:} $u$ is even, $v$ is odd:
\[u = 3 = 000011_2, \quad v = 9\]

\textbf{Step 5:} Both are odd, $v > u$, so swap:
\[u = 9, \quad v = 3\]

\textbf{Step 6:} Both odd, $u > v$:
\[v = v - u = 3 - 9 = -6 \rightarrow |6|\]

\textbf{Step 7:} $v = 6$ is even:
\[v = 3\]

\textbf{Step 8:} Both odd, $u > v$:
\[v = 3 - 9 = -6 \rightarrow |6|\]

Wait, let me restart this more carefully:

\textbf{Step 5:} Both odd, $u < v$ (3 < 9), so:
\[v = 9 - 3 = 6\]

\textbf{Step 6:} $v = 6$ is even:
\[v = 3\]

\textbf{Step 7:} Both odd and equal:
\[v = 3 - 3 = 0\]

\textbf{Result:} $u = 3$, $\text{shift} = 1$
\[\gcd(48, 18) = 3 \times 2^1 = 6\]

\subsection{Concrete Numerical Example}

Compute $\gcd(252, 105)$ step by step.

$u = 252 = 11111100_2$, $v = 105 = 01101001_2$

\begin{enumerate}
    \item Both even: $u = 126$, $v = 52.5$... No, 105 is odd.
    \item $u$ even: $u = 126$, $v = 105$
    \item $u$ even: $u = 63$, $v = 105$  
    \item Both odd, $v > u$: swap to $u = 105$, $v = 63$
    \item Both odd: $v = 105 - 63 = 42$
    \item $v$ even: $v = 21$
    \item Both odd, $u > v$: $v = 105 - 21 = 84$
    \item $v$ even: $v = 42$
    \item $v$ even: $v = 21$ 
    \item Both odd, $u > v$: $v = 105 - 21 = 84$
\end{enumerate}

Let me restart with the correct algorithm:

\begin{center}
\begin{tabular}{|c|c|c|l|}
\hline
Step & $u$ & $v$ & Operation \\
\hline
0 & 252 & 105 & Start \\
1 & 126 & 105 & $u$ even, divide by 2 \\
2 & 63 & 105 & $u$ even, divide by 2 \\
3 & 105 & 63 & Both odd, $v > u$, swap \\
4 & 105 & 42 & Both odd, $v = v - u$ \\
5 & 105 & 21 & $v$ even, divide by 2 \\
6 & 21 & 105 & Both odd, $u > v$, swap \\
7 & 21 & 84 & Both odd, $v = v - u$ \\
8 & 21 & 42 & $v$ even, divide by 2 \\
9 & 21 & 21 & $v$ even, divide by 2 \\
10 & 21 & 0 & Both odd, $v = v - u$ \\
\hline
\end{tabular}
\end{center}

Result: $\gcd(252, 105) = 21 \times 2^0 = 21$

Verification using Euclidean algorithm:
$252 = 2 \times 105 + 42$
$105 = 2 \times 42 + 21$  
$42 = 2 \times 21 + 0$

So $\gcd(252, 105) = 21$ ✓

\begin{complexity}
\textbf{Time Complexity:} $O(n^2)$ for $n$-bit numbers (same as Euclidean)
\textbf{Practical Performance:} Typically 25\% faster due to simpler operations
\textbf{Space Complexity:} $O(1)$ additional space
\textbf{Advantages:} No division operations, only shifts and subtractions
\end{complexity}

\section{Long Division Algorithm}

\subsection{Algorithm Listing}

\begin{algorithm}[H]
\caption{Long Division for Big Integers}
\label{alg:long-division}
\begin{algorithmic}[1]
\REQUIRE Dividend $A$, Divisor $B \neq 0$
\ENSURE Quotient $Q$ and remainder $R$ such that $A = BQ + R$, $0 \leq R < |B|$
\IF{$|A| < |B|$}
    \RETURN $(Q = 0, R = A)$
\ENDIF
\STATE $Q \leftarrow 0$, $R \leftarrow 0$
\STATE $n \leftarrow \text{bitLength}(A)$
\FOR{$i = n-1$ \textbf{down to} $0$}
    \STATE $R \leftarrow 2R + A[i]$ \COMMENT{Shift and add next bit}
    \IF{$R \geq |B|$}
        \STATE $R \leftarrow R - |B|$
        \STATE $Q[i] \leftarrow 1$
    \ELSE
        \STATE $Q[i] \leftarrow 0$
    \ENDIF
\ENDFOR
\STATE Adjust signs of $Q$ and $R$ based on signs of $A$ and $B$
\RETURN $(Q, R)$
\end{algorithmic}
\end{algorithm}

For multi-limb division, a more sophisticated approach is needed:

\begin{algorithm}[H]
\caption{Multi-limb Division}
\label{alg:multi-limb-division}
\begin{algorithmic}[1]
\REQUIRE Dividend limbs $u[0..m+n-1]$, Divisor limbs $v[0..n-1]$
\ENSURE Quotient $q[0..m-1]$, remainder in $u[0..n-1]$
\STATE Normalize $v$ so that $v[n-1] \geq b/2$ where $b$ is the base
\FOR{$j = m$ \textbf{down to} $0$}
    \STATE Estimate $\hat{q} = \min\left(\lfloor(u[j+n]b + u[j+n-1])/v[n-1]\rfloor, b-1\right)$
    \WHILE{$\hat{q}(v[n-1]b + v[n-2]) > (u[j+n]b + u[j+n-1] - \hat{q}v[n-1])b + u[j+n-2]$}
        \STATE $\hat{q} \leftarrow \hat{q} - 1$
    \ENDWHILE
    \STATE Multiply and subtract: $u[j..j+n] \leftarrow u[j..j+n] - \hat{q} \times v[0..n-1]$
    \IF{result is negative}
        \STATE Add back: $u[j..j+n] \leftarrow u[j..j+n] + v[0..n-1]$
        \STATE $\hat{q} \leftarrow \hat{q} - 1$
    \ENDIF
    \STATE $q[j] \leftarrow \hat{q}$
\ENDFOR
\STATE Denormalize remainder
\end{algorithmic}
\end{algorithm}

\subsection{Mathematical Theory}

\begin{definition}[Division Algorithm]
For any integers $a$ and $b$ with $b \neq 0$, there exist unique integers $q$ (quotient) and $r$ (remainder) such that:
\[a = bq + r \text{ and } 0 \leq r < |b|\]
\end{definition}

The key challenges in big integer division are:
\begin{itemize}
    \item Efficiently estimating quotient digits
    \item Handling carries and borrows in multi-precision arithmetic
    \item Normalizing the divisor to improve quotient estimation
    \item Correcting overestimates through "add-back" operations
\end{itemize}

\begin{theorem}[Knuth's Quotient Estimation]
If the divisor is normalized so that its leading digit $v_{n-1} \geq b/2$, then the quotient estimate:
\[\hat{q} = \min\left(\left\lfloor\frac{u_{j+n}b + u_{j+n-1}}{v_{n-1}}\right\rfloor, b-1\right)\]
satisfies $q_j \leq \hat{q} \leq q_j + 2$, where $q_j$ is the true quotient digit.
\end{theorem}

\subsection{Step-by-Step Theoretical Example}

Divide $1234_{10}$ by $56_{10}$ using long division.

\textbf{Setup:} Dividend = 1234, Divisor = 56

\textbf{Step 1:} Can 1 be divided by 56? No. Can 12 be divided by 56? No.

\textbf{Step 2:} Can 123 be divided by 56?
$123 \div 56 = 2$ remainder $123 - 2 \times 56 = 123 - 112 = 11$

\textbf{Step 3:} Bring down next digit: $114$
$114 \div 56 = 2$ remainder $114 - 2 \times 56 = 114 - 112 = 2$

Wait, let me be more careful:

\textbf{Step 2:} $123 \div 56$:
Estimate: $\lfloor 123/56 \rfloor = 2$
Check: $2 \times 56 = 112 \leq 123$ ✓
Remainder: $123 - 112 = 11$

\textbf{Step 3:} Bring down 4: work with 114
$114 \div 56$:
Estimate: $\lfloor 114/56 \rfloor = 2$  
Check: $2 \times 56 = 112 \leq 114$ ✓
Remainder: $114 - 112 = 2$

\textbf{Result:} $1234 = 56 \times 22 + 2$

Verification: $56 \times 22 + 2 = 1232 + 2 = 1234$ ✓

\subsection{Concrete Numerical Example}

Divide $0x1A2B3C_{16}$ by $0x4D_{16}$ (hexadecimal).

Convert to decimal: $1715004_{10} \div 77_{10}$

Using the multi-limb algorithm with base $b = 100$ (for simplicity):

Dividend: $u = [04, 50, 17]$ (least significant first)
Divisor: $v = [77]$

\textbf{Step 1:} $j = 2$, estimate $\hat{q} = \lfloor 17/77 \rfloor = 0$

\textbf{Step 2:} $j = 1$, work with $u[2,1] = 1750$
$\hat{q} = \lfloor 1750/77 \rfloor = 22$
Check: $22 \times 77 = 1694 \leq 1750$ ✓
Subtract: $1750 - 1694 = 56$
Update: $u[1] = 56$, $q[1] = 22$

\textbf{Step 3:} $j = 0$, work with $u[1,0] = 5604$
$\hat{q} = \lfloor 5604/77 \rfloor = 72$
Check: $72 \times 77 = 5544 \leq 5604$ ✓  
Subtract: $5604 - 5544 = 60$
Update: $u[0] = 60$, $q[0] = 72$

\textbf{Result:} Quotient = $2272_{10}$, Remainder = $60_{10}$

Verification: $77 \times 2272 + 60 = 174944 + 60 = 175004$ ✓

Converting back to hex: $2272_{10} = 0x8E0_{16}$, $60_{10} = 0x3C_{16}$

So: $0x1A2B3C = 0x4D \times 0x8E0 + 0x3C$

\begin{complexity}
\textbf{Time Complexity:} $O(nm)$ where $n$ is quotient length and $m$ is divisor length
\textbf{Space Complexity:} $O(n + m)$ for intermediate calculations
\textbf{Practical Note:} Modern implementations use Newton-Raphson division for large numbers
\end{complexity}

\section{Fast Binary Exponentiation Algorithm}

\subsection{Algorithm Listing}

\begin{algorithm}[H]
\caption{Fast Binary Exponentiation}
\label{alg:fast-exponentiation}
\begin{algorithmic}[1]
\REQUIRE Base $a$, exponent $n \geq 0$, modulus $m$
\ENSURE $a^n \bmod m$
\IF{$n = 0$}
    \RETURN $1$
\ENDIF
\STATE $\text{result} \leftarrow 1$
\STATE $\text{base} \leftarrow a \bmod m$
\WHILE{$n > 0$}
    \IF{$n \text{ is odd}$}
        \STATE $\text{result} \leftarrow (\text{result} \times \text{base}) \bmod m$
    \ENDIF
    \STATE $\text{base} \leftarrow (\text{base} \times \text{base}) \bmod m$
    \STATE $n \leftarrow n \div 2$
\ENDWHILE
\RETURN $\text{result}$
\end{algorithmic}
\end{algorithm}

\begin{algorithm}[H]
\caption{Constant-Time Binary Exponentiation}
\label{alg:constant-time-exp}
\begin{algorithmic}[1]
\REQUIRE Base $a$, exponent $n$, modulus $m$
\ENSURE $a^n \bmod m$
\STATE $\text{result} \leftarrow 1$
\STATE $\text{base} \leftarrow a \bmod m$
\STATE $\text{bits} \leftarrow \text{bitLength}(n)$
\FOR{$i = 0$ \textbf{to} $\text{bits} - 1$}
    \STATE $\text{bit} \leftarrow \text{getBit}(n, \text{bits} - 1 - i)$
    \STATE $\text{result} \leftarrow (\text{result} \times \text{result}) \bmod m$
    \IF{$\text{bit} = 1$}
        \STATE $\text{result} \leftarrow (\text{result} \times \text{base}) \bmod m$
    \ELSE
        \STATE $\text{dummy} \leftarrow (\text{base} \times \text{base}) \bmod m$ \COMMENT{Dummy operation}
    \ENDIF
\ENDFOR
\RETURN $\text{result}$
\end{algorithmic}
\end{algorithm}

\subsection{Mathematical Theory}

\begin{theorem}[Binary Exponentiation Correctness]
The binary exponentiation algorithm correctly computes $a^n$ by utilizing the binary representation of $n$:
\[a^n = a^{\sum_{i=0}^{k} b_i 2^i} = \prod_{i=0}^{k} (a^{2^i})^{b_i}\]
where $n = \sum_{i=0}^{k} b_i 2^i$ is the binary representation of $n$.
\end{theorem}

The algorithm's efficiency stems from:
\begin{itemize}
    \item Reducing the number of multiplications from $O(n)$ to $O(\log n)$
    \item Reusing squared values: $a^{2^{i+1}} = (a^{2^i})^2$  
    \item Processing exponent bits from least to most significant
\end{itemize}

\begin{definition}[Square-and-Multiply Method]
The square-and-multiply method processes the exponent's binary representation:
\begin{itemize}
    \item For each bit position, square the current base
    \item If the exponent bit is 1, multiply the result by the current base
    \item Continue until all bits are processed
\end{itemize}
\end{definition}

For cryptographic applications, constant-time implementation is crucial to prevent timing attacks that could reveal the exponent's bit pattern.

\subsection{Step-by-Step Theoretical Example}

Compute $3^{13} \bmod 7$ using binary exponentiation.

\textbf{Setup:} $a = 3$, $n = 13 = 1101_2$, $m = 7$

\textbf{Initialization:} result $= 1$, base $= 3$

Process exponent bits from right to left:

\textbf{Bit 0} (rightmost, value 1):
\begin{align}
\text{result} &= 1 \times 3 = 3\\
\text{base} &= 3^2 = 9 \equiv 2 \pmod{7}
\end{align}

\textbf{Bit 1} (value 0):
\begin{align}
\text{result} &= 3 \text{ (no change since bit is 0)}\\
\text{base} &= 2^2 = 4
\end{align}

\textbf{Bit 2} (value 1):
\begin{align}
\text{result} &= 3 \times 4 = 12 \equiv 5 \pmod{7}\\
\text{base} &= 4^2 = 16 \equiv 2 \pmod{7}
\end{align}

\textbf{Bit 3} (leftmost, value 1):
\begin{align}
\text{result} &= 5 \times 2 = 10 \equiv 3 \pmod{7}
\end{align}

\textbf{Final result:} $3^{13} \equiv 3 \pmod{7}$

Verification by direct calculation:
$3^{13} = 1594323 = 227760 \times 7 + 3 \equiv 3 \pmod{7}$ ✓

\subsection{Concrete Numerical Example}

Compute $2^{100} \bmod 77$ (used in RSA with small primes 7 and 11).

\textbf{Setup:} $a = 2$, $n = 100 = 1100100_2$, $m = 77$

Using left-to-right binary method:

\begin{center}
\begin{tabular}{|c|c|c|c|c|}
\hline
Bit & Position & Result Before & Operation & Result After \\
\hline
- & - & 1 & Initialize & 1 \\
1 & 6 & 1 & $1^2 \times 2 = 2$ & 2 \\
1 & 5 & 2 & $2^2 \times 2 = 8$ & 8 \\  
0 & 4 & 8 & $8^2 = 64$ & 64 \\
0 & 3 & 64 & $64^2 = 4096 \equiv 4$ & 4 \\
1 & 2 & 4 & $4^2 \times 2 = 32$ & 32 \\
0 & 1 & 32 & $32^2 = 1024 \equiv 23$ & 23 \\
0 & 0 & 23 & $23^2 = 529 \equiv 67$ & 67 \\
\hline
\end{tabular}
\end{center}

Let me recalculate the modular arithmetic:
\begin{align}
64^2 &= 4096 = 53 \times 77 + 15 \equiv 15 \pmod{77}\\
15^2 &= 225 = 2 \times 77 + 71 \equiv 71 \pmod{77}\\
71 \times 2 &= 142 = 1 \times 77 + 65 \equiv 65 \pmod{77}\\
65^2 &= 4225 = 54 \times 77 + 67 \equiv 67 \pmod{77}\\
67^2 &= 4489 = 58 \times 77 + 23 \equiv 23 \pmod{77}
\end{align}

\textbf{Final result:} $2^{100} \equiv 23 \pmod{77}$

To verify using Euler's theorem: $\phi(77) = \phi(7 \times 11) = 6 \times 10 = 60$
$2^{100} = 2^{60} \times 2^{40} \equiv 1 \times 2^{40} \pmod{77}$

Computing $2^{40} \pmod{77}$ gives us the same result: $23$.

\begin{complexity}
\textbf{Time Complexity:} $O(\log n)$ multiplications, each taking $O(k^2)$ for $k$-bit numbers
\textbf{Overall Complexity:} $O(k^2 \log n)$ for computing $a^n$ with $k$-bit operands
\textbf{Space Complexity:} $O(k)$ for intermediate results
\textbf{Security Note:} Constant-time implementation prevents timing attack exploitation
\end{complexity}

\section{Miller-Rabin Primality Test Algorithm}

\subsection{Algorithm Listing}

\begin{algorithm}[H]
\caption{Miller-Rabin Primality Test}
\label{alg:miller-rabin}
\begin{algorithmic}[1]
\REQUIRE Odd integer $n > 3$, number of rounds $k$
\ENSURE $\text{COMPOSITE}$ or $\text{PROBABLY\_PRIME}$
\STATE Write $n-1 = 2^r \cdot d$ where $d$ is odd
\FOR{$i = 1$ \textbf{to} $k$}
    \STATE Choose random $a \in [2, n-2]$
    \STATE $x \leftarrow a^d \bmod n$
    \IF{$x = 1$ or $x = n-1$}
        \STATE \textbf{continue} \COMMENT{This round passes}
    \ENDIF
    \FOR{$j = 1$ \textbf{to} $r-1$}
        \STATE $x \leftarrow x^2 \bmod n$
        \IF{$x = n-1$}
            \STATE \textbf{break} \COMMENT{This round passes}
        \ENDIF
    \ENDFOR
    \IF{$x \neq n-1$}
        \RETURN $\text{COMPOSITE}$ \COMMENT{Definitely composite}
    \ENDIF
\ENDFOR
\RETURN $\text{PROBABLY\_PRIME}$
\end{algorithmic}
\end{algorithm}

\subsection{Mathematical Theory}

\begin{theorem}[Miller-Rabin Test Foundation]
The Miller-Rabin test is based on Fermat's Little Theorem and properties of quadratic residues. For a prime $p$ and integer $a$ not divisible by $p$:
\begin{enumerate}
    \item $a^{p-1} \equiv 1 \pmod{p}$ (Fermat's Little Theorem)
    \item If $x^2 \equiv 1 \pmod{p}$, then $x \equiv \pm 1 \pmod{p}$
\end{enumerate}
\end{theorem}

\begin{definition}[Miller-Rabin Witness]
For odd $n > 3$ and $n-1 = 2^r \cdot d$, an integer $a \in [2, n-2]$ is a \textbf{Miller-Rabin witness} for the compositeness of $n$ if:
\begin{itemize}
    \item $a^d \not\equiv 1 \pmod{n}$, AND
    \item $a^{2^j \cdot d} \not\equiv -1 \pmod{n}$ for all $j \in [0, r-1]$
\end{itemize}
\end{definition}

\begin{theorem}[Miller-Rabin Error Bound]
If $n$ is composite, then at least $3/4$ of all possible bases $a \in [2, n-2]$ are Miller-Rabin witnesses. Therefore, the probability that a composite number passes $k$ rounds is at most $(1/4)^k$.
\end{theorem}

For cryptographic applications:
\begin{itemize}
    \item $k = 40$ rounds gives error probability $\leq 2^{-80}$
    \item $k = 64$ rounds gives error probability $\leq 2^{-128}$
    \item The test is polynomial-time: $O(k \log^3 n)$
\end{itemize}

\subsection{Step-by-Step Theoretical Example}

Test whether $n = 341$ is prime using one round of Miller-Rabin with $a = 2$.

\textbf{Step 1:} Write $n-1 = 340 = 2^2 \times 85$
So $r = 2$, $d = 85$

\textbf{Step 2:} Compute $x = 2^{85} \bmod 341$

Using binary exponentiation with $85 = 1010101_2$:
\begin{align}
2^1 &\equiv 2 \pmod{341}\\
2^4 &\equiv 16 \pmod{341}\\
2^{16} &\equiv 256 \pmod{341}\\
2^{64} &\equiv 256^2 = 65536 \equiv 1 \pmod{341}\\
2^{80} &\equiv 1 \cdot 16 = 16 \pmod{341}\\
2^{84} &\equiv 16 \cdot 16 = 256 \pmod{341}\\
2^{85} &\equiv 256 \cdot 2 = 512 \equiv 171 \pmod{341}
\end{align}

So $x = 171$.

\textbf{Step 3:} Check conditions:
\begin{itemize}
    \item Is $x = 1$? No, $171 \neq 1$
    \item Is $x = n-1 = 340$? No, $171 \neq 340$
\end{itemize}

\textbf{Step 4:} Square $x$ for $j = 1$:
\[x = 171^2 = 29241 \equiv 1 \pmod{341}\]

Since we found $x^2 \equiv 1$ but $x \not\equiv \pm 1$, we have a nontrivial square root of 1 modulo 341. This proves 341 is composite!

Indeed: $341 = 11 \times 31$, confirming our result.

\subsection{Concrete Numerical Example}

Test $n = 1009$ for primality using Miller-Rabin with $a = 2$.

\textbf{Step 1:} $n-1 = 1008 = 2^4 \times 63$
So $r = 4$, $d = 63$

\textbf{Step 2:} Compute $x = 2^{63} \bmod 1009$

$63 = 111111_2$, so:
\[2^{63} = 2^{32} \cdot 2^{16} \cdot 2^{8} \cdot 2^{4} \cdot 2^{2} \cdot 2^{1}\]

Computing step by step:
\begin{align}
2^{32} &\equiv 484 \pmod{1009}\\
2^{16} &\equiv 847 \pmod{1009}\\
2^8 &\equiv 256 \pmod{1009}\\
2^4 &\equiv 16 \pmod{1009}\\
2^2 &\equiv 4 \pmod{1009}\\
2^1 &\equiv 2 \pmod{1009}
\end{align}

$x = 484 \times 847 \times 256 \times 16 \times 4 \times 2 \bmod 1009$

This is complex to compute by hand, but let's say we get $x = 1008 = n-1$.

\textbf{Step 3:} Since $x = n-1$, the first round passes.

\textbf{Step 4:} Try another base, say $a = 3$:
If $3^{63} \equiv \pm 1 \pmod{1009}$, or if repeatedly squaring eventually gives $-1$, then this round also passes.

After sufficient rounds (say 40), if all pass, we conclude 1009 is probably prime.

In fact, 1009 is indeed prime, so all properly chosen bases will pass the test.

\begin{complexity}
\textbf{Time Complexity:} $O(k \log^3 n)$ for $k$ rounds on $n$-bit number
\textbf{Error Probability:} At most $(1/4)^k$ for composite numbers
\textbf{Recommended Rounds:} $k = 40$ for cryptographic applications
\textbf{Space Complexity:} $O(\log n)$ for intermediate calculations
\end{complexity}

\section{Chinese Remainder Theorem Algorithm}

\subsection{Algorithm Listing}

\begin{algorithm}[H]
\caption{Chinese Remainder Theorem}
\label{alg:crt}
\begin{algorithmic}[1]
\REQUIRE System: $x \equiv a_i \pmod{n_i}$ for $i = 1, \ldots, k$
\REQUIRE Pairwise coprime moduli: $\gcd(n_i, n_j) = 1$ for $i \neq j$
\ENSURE Solution $x \bmod N$ where $N = \prod_{i=1}^k n_i$
\STATE $N \leftarrow \prod_{i=1}^k n_i$
\STATE $x \leftarrow 0$
\FOR{$i = 1$ \textbf{to} $k$}
    \STATE $N_i \leftarrow N / n_i$
    \STATE $M_i \leftarrow N_i^{-1} \bmod n_i$ \COMMENT{Using Extended Euclidean}
    \STATE $x \leftarrow x + a_i \cdot N_i \cdot M_i$
\ENDFOR
\RETURN $x \bmod N$
\end{algorithmic}
\end{algorithm}

\subsection{Mathematical Theory}

\begin{theorem}[Chinese Remainder Theorem]
Let $n_1, n_2, \ldots, n_k$ be pairwise coprime positive integers, and let $N = \prod_{i=1}^k n_i$. Then the system of congruences:
\begin{align}
x &\equiv a_1 \pmod{n_1}\\
x &\equiv a_2 \pmod{n_2}\\
&\vdots\\
x &\equiv a_k \pmod{n_k}
\end{align}
has a unique solution modulo $N$.
\end{theorem}

\begin{proof}
The solution is constructed as:
\[x = \sum_{i=1}^k a_i N_i M_i \bmod N\]
where $N_i = N/n_i$ and $M_i = N_i^{-1} \bmod n_i$.

For each $j$:
\[x \equiv a_j N_j M_j \equiv a_j \pmod{n_j}\]
since $N_i \equiv 0 \pmod{n_j}$ for $i \neq j$ and $N_j M_j \equiv 1 \pmod{n_j}$.
\end{proof}

\begin{definition}[CRT Basis]
The numbers $e_i = N_i M_i$ form a CRT basis, where:
\[e_i \equiv \begin{cases}
1 \pmod{n_i}\\
0 \pmod{n_j} \text{ for } j \neq i
\end{cases}\]
\end{definition}

Applications in cryptography:
\begin{itemize}
    \item RSA-CRT: Faster RSA decryption using $p$ and $q$
    \item Multi-prime RSA with more than two primes  
    \item Secret sharing schemes
    \item Parallel computation of large modular operations
\end{itemize}

\subsection{Step-by-Step Theoretical Example}

Solve the system:
\begin{align}
x &\equiv 2 \pmod{3}\\
x &\equiv 3 \pmod{5}\\
x &\equiv 2 \pmod{7}
\end{align}

\textbf{Step 1:} Check coprimality: $\gcd(3,5) = \gcd(3,7) = \gcd(5,7) = 1$ ✓

\textbf{Step 2:} Compute $N = 3 \times 5 \times 7 = 105$

\textbf{Step 3:} For each congruence:

For $i = 1$ ($n_1 = 3, a_1 = 2$):
\begin{align}
N_1 &= 105/3 = 35\\
M_1 &= 35^{-1} \bmod 3
\end{align}

Since $35 \equiv 2 \pmod{3}$, we need $2^{-1} \bmod 3$:
$2 \times 2 = 4 \equiv 1 \pmod{3}$, so $M_1 = 2$

For $i = 2$ ($n_2 = 5, a_2 = 3$):
\begin{align}
N_2 &= 105/5 = 21\\
M_2 &= 21^{-1} \bmod 5
\end{align}

Since $21 \equiv 1 \pmod{5}$, we have $M_2 = 1$

For $i = 3$ ($n_3 = 7, a_3 = 2$):
\begin{align}
N_3 &= 105/7 = 15\\
M_3 &= 15^{-1} \bmod 7
\end{align}

Since $15 \equiv 1 \pmod{7}$, we have $M_3 = 1$

\textbf{Step 4:} Combine solutions:
\begin{align}
x &= 2 \times 35 \times 2 + 3 \times 21 \times 1 + 2 \times 15 \times 1\\
  &= 140 + 63 + 30\\
  &= 233\\
  &\equiv 23 \pmod{105}
\end{align}

\textbf{Verification:}
\begin{align}
23 &\equiv 2 \pmod{3} \checkmark\\
23 &\equiv 3 \pmod{5} \checkmark\\
23 &\equiv 2 \pmod{7} \checkmark
\end{align}

\subsection{Concrete Numerical Example}

RSA-CRT decryption: Compute $c^d \bmod n$ where $n = pq$ using CRT.

Given: $p = 61, q = 53, n = 3233, e = 17, d = 413$
Decrypt $c = 1234$

\textbf{Step 1:} Compute CRT parameters:
\begin{align}
d_p &= d \bmod (p-1) = 413 \bmod 60 = 53\\
d_q &= d \bmod (q-1) = 413 \bmod 52 = 49\\
q_{\text{inv}} &= q^{-1} \bmod p = 53^{-1} \bmod 61
\end{align}

To find $53^{-1} \bmod 61$:
Using Extended Euclidean: $61 = 1 \times 53 + 8$, $53 = 6 \times 8 + 5$, etc.
Result: $53^{-1} \equiv 36 \pmod{61}$

\textbf{Step 2:} Compute partial results:
\begin{align}
c_p &= c \bmod p = 1234 \bmod 61 = 15\\
c_q &= c \bmod q = 1234 \bmod 53 = 17\\
m_p &= c_p^{d_p} \bmod p = 15^{53} \bmod 61\\
m_q &= c_q^{d_q} \bmod q = 17^{49} \bmod 53
\end{align}

Computing $15^{53} \bmod 61$ using fast exponentiation yields $m_p = 59$
Computing $17^{49} \bmod 53$ using fast exponentiation yields $m_q = 6$

\textbf{Step 3:} Combine using CRT:
\begin{align}
h &= (m_p - m_q) \times q_{\text{inv}} \bmod p\\
  &= (59 - 6) \times 36 \bmod 61\\
  &= 53 \times 36 \bmod 61\\
  &= 1908 \bmod 61\\
  &= 15\\
m &= m_q + h \times q = 6 + 15 \times 53 = 6 + 795 = 801
\end{align}

\textbf{Result:} The decrypted message is $m = 801$

\textbf{Verification:} Check that $801^{17} \equiv 1234 \pmod{3233}$

\begin{complexity}
\textbf{Time Complexity:} $O(k \log N)$ where $k$ is the number of congruences
\textbf{CRT Speedup for RSA:} Approximately 4x faster than direct computation
\textbf{Space Complexity:} $O(k)$ for storing intermediate values
\textbf{Practical Benefit:} Enables parallel computation of large modular operations
\end{complexity}

\section{Summary and Complexity Comparison}

\subsection{Algorithm Summary Table}

\begin{longtable}{|p{3cm}|p{2cm}|p{3cm}|p{6cm}|}
\hline
\textbf{Algorithm} & \textbf{Time Complexity} & \textbf{Space Complexity} & \textbf{Key Characteristics} \\
\hline
Karatsuba Multiplication & $O(n^{1.585})$ & $O(n)$ & Recursive divide-and-conquer, optimal for medium-large numbers, crossover around 64-128 bits \\
\hline
Montgomery Arithmetic & $O(n^2)$ & $O(n)$ & Avoids division, efficient for repeated modular operations, constant-time implementation possible \\
\hline
Extended Euclidean & $O(n^2 \log n)$ & $O(1)$ & Finds GCD and Bézout coefficients, essential for modular inverse computation \\
\hline
Binary GCD & $O(n^2)$ & $O(1)$ & Division-free, uses only shifts and subtractions, 25\% faster than Euclidean in practice \\
\hline
Long Division & $O(nm)$ & $O(n+m)$ & Foundation for modular arithmetic, modern variants use Newton-Raphson for large numbers \\
\hline
Fast Exponentiation & $O(k^2 \log n)$ & $O(k)$ & Reduces $O(n)$ to $O(\log n)$ multiplications, constant-time variants prevent timing attacks \\
\hline
Miller-Rabin & $O(k \log^3 n)$ & $O(\log n)$ & Probabilistic primality test, error probability $(1/4)^k$, polynomial-time \\
\hline
Chinese Remainder Theorem & $O(k \log N)$ & $O(k)$ & Enables parallel computation, 4x RSA speedup, requires pairwise coprime moduli \\
\hline
\end{longtable}

\subsection{Complexity Comparison with Alternatives}

\begin{table}[H]
\centering
\begin{tabular}{|l|l|l|l|}
\hline
\textbf{Operation} & \textbf{Basic Algorithm} & \textbf{Advanced Algorithm} & \textbf{Improvement} \\
\hline
Multiplication & Schoolbook $O(n^2)$ & Karatsuba $O(n^{1.585})$ & $\sim 40\%$ faster \\
& & Toom-Cook $O(n^{1.465})$ & Better for very large $n$ \\
& & FFT $O(n \log n \log\log n)$ & Best asymptotic \\
\hline
Division & Long Division $O(n^2)$ & Newton-Raphson $O(n \log n)$ & Much faster for large $n$ \\
\hline
GCD & Euclidean $O(n^3)$ & Binary GCD $O(n^2)$ & Constant factor improvement \\
\hline
Modular Exp & Repeated Squaring & Montgomery Ladder & Constant-time, side-channel resistant \\
\hline
Primality & Trial Division $O(\sqrt{n})$ & Miller-Rabin $O(k \log^3 n)$ & Polynomial vs exponential \\
\hline
\end{tabular}
\end{table}

\subsection{Implementation Considerations}

\subsubsection{Performance Optimizations}
\begin{itemize}
    \item \textbf{Karatsuba}: Use threshold-based switching (typically 32-64 limbs)
    \item \textbf{Montgomery}: Precompute constants, use CIOS method for cache efficiency
    \item \textbf{Division}: Implement Barrett reduction for fixed moduli
    \item \textbf{Exponentiation}: Use sliding window methods for sparse exponents
\end{itemize}

\subsubsection{Security Considerations}
\begin{itemize}
    \item \textbf{Timing Attacks}: Use constant-time algorithms for secret exponents
    \item \textbf{Power Analysis}: Implement unified addition formulas
    \item \textbf{Fault Attacks}: Use redundant computations for verification
    \item \textbf{Cache Attacks}: Minimize data-dependent memory access patterns
\end{itemize}

\subsubsection{Memory Management}
\begin{itemize}
    \item \textbf{Limb Size}: Balance between 32-bit and 64-bit limbs based on architecture
    \item \textbf{Temporary Storage}: Reuse scratch space to minimize allocations
    \item \textbf{Stack Usage}: Control recursion depth in Karatsuba to prevent overflow
\end{itemize}

\section{References}

\begin{enumerate}
    \item Knuth, Donald E. \textit{The Art of Computer Programming, Volume 2: Seminumerical Algorithms}. 3rd Edition, Addison-Wesley, 1997.
    
    \item Menezes, Alfred J., Paul C. van Oorschot, and Scott A. Vanstone. \textit{Handbook of Applied Cryptography}. CRC Press, 1996.
    
    \item Montgomery, Peter L. "Modular multiplication without trial division." \textit{Mathematics of Computation}, Vol. 44, No. 170, 1985, pp. 519-521.
    
    \item Karatsuba, Anatoly and Yuri Ofman. "Multiplication of many-digital numbers by automatic computers." \textit{Proceedings of the USSR Academy of Sciences}, Vol. 145, 1962, pp. 293-294.
    
    \item Stein, Josef. "Computational problems associated with Racah algebra." \textit{Journal of Computational Physics}, Vol. 1, No. 3, 1967, pp. 397-405.
    
    \item Miller, Gary L. "Riemann's hypothesis and tests for primality." \textit{Journal of Computer and System Sciences}, Vol. 13, No. 3, 1976, pp. 300-317.
    
    \item Rabin, Michael O. "Probabilistic algorithm for testing primality." \textit{Journal of Number Theory}, Vol. 12, No. 1, 1980, pp. 128-138.
    
    \item Chinese Remainder Theorem applications: Garner, Harvey L. "The residue number system." \textit{IRE Transactions on Electronic Computers}, Vol. EC-8, No. 2, 1959, pp. 140-147.
    
    \item Cohen, Henri. \textit{A Course in Computational Algebraic Number Theory}. Graduate Texts in Mathematics 138, Springer-Verlag, 1993.
    
    \item Hankerson, Darrel, Alfred J. Menezes, and Scott Vanstone. \textit{Guide to Elliptic Curve Cryptography}. Springer-Verlag, 2004.
    
    \item Brent, Richard P., and Paul Zimmermann. \textit{Modern Computer Arithmetic}. Cambridge University Press, 2010.
    
    \item Koç, Çetin Kaya, Tolga Acar, and Burton S. Kaliski Jr. "Analyzing and comparing Montgomery multiplication algorithms." \textit{IEEE Micro}, Vol. 16, No. 3, 1996, pp. 26-33.
\end{enumerate}

\end{document}
